% ---
% title: Picking bonds with and without replacement
% tags: finance, hypergeometric, binomial
% ---

\begin{problem}
	\begin{context}
		Rating agencies grade bonds by how likely the issuer is to pay back.
		Bonds with a low credit rating are called \emph{high-yield} or
		\emph{junk} bonds: their issuers are more likely to default, so these
		bonds must pay a higher interest rate to attract investors.
	\end{context}
	A fund holds 50 bonds, of which 10 are junk bonds. An analyst picks 4 of
	the 50 bonds at random to review. Let $X$ be the number of junk bonds
	among them.
	(a) The analyst picks 4 different bonds (without replacement). What is the
	distribution of $X$? Find $P(X=0)$ and $P(X=1)$.
	(b) The analyst picks one bond at random 4 times, and the same bond may be
	picked again (with replacement): the 4 draws are independent, and each is
	uniform over all 50 bonds. Now let $X$ be the number of the 4 draws that
	select a junk bond, counting repeated selections separately. What is the
	distribution of $X$ now? Find $P(X=0)$ and $P(X=1)$.
	(c) Compare the results of (a) and (b). Find $\E[X]$ in both cases.
\end{problem}

\begin{solution}
	\idea Sampling \emph{without} replacement from a finite group gives a
	hypergeometric distribution. Sampling \emph{with} replacement gives
	independent trials with the same success probability, so a binomial
	distribution.

	\step{(a)} The 4 bonds are different: the analyst draws a sample of $n=4$
	from $N=50$ bonds, of which $m=10$ are ``special'' (junk). The number of
	special items in such a sample is hypergeometric with parameters
	$N=50$, $m=10$, $n=4$:
	\[ P(X=i)=\frac{\dbinom{10}{i}\dbinom{40}{4-i}}{\dbinom{50}{4}},
		\qquad i=0,1,2,3,4. \]
	Why: every set of 4 bonds is equally likely, and there are
	$\binom{50}{4}=230300$ such sets. To get exactly $i$ junk bonds, choose
	$i$ of the 10 junk bonds and $4-i$ of the 40 other bonds.

	For $i=0$ all 4 bonds come from the 40 good ones, $\binom{40}{4}=91390$:
	\[ P(X=0)=\frac{\dbinom{10}{0}\dbinom{40}{4}}{\dbinom{50}{4}}
		=\frac{91390}{230300}\approx0.3968. \]
	For $i=1$ choose 1 of the 10 junk bonds and 3 of the 40 others,
	$\binom{40}{3}=9880$:
	\[ P(X=1)=\frac{\dbinom{10}{1}\dbinom{40}{3}}{\dbinom{50}{4}}
		=\frac{10\cdot9880}{230300}=\frac{98800}{230300}\approx0.4290. \]

	\step{(b)} Now each pick is made uniformly from all 50 bonds, whatever
	happened before, and $X$ counts the draws that select a junk bond (a junk
	bond picked twice counts twice). So the 4 picks are independent trials, and in each trial the
	chance of a junk bond is the same, $p=\frac{10}{50}=0.2$. The number of
	``successes'' (junk bonds) in $n=4$ independent trials with success
	probability $p=0.2$ is binomial with parameters $(4,\,0.2)$:
	\[ P(X=i)=\binom{4}{i}0.2^i\,0.8^{4-i},\qquad i=0,1,2,3,4. \]
	Plug in:
	\[ P(X=0)=0.8^4=0.4096, \]
	\[ P(X=1)=\binom41\cdot0.2\cdot0.8^3=4\cdot0.2\cdot0.512=0.4096. \]
	(It is a coincidence of these numbers that the two are equal.)

	\step{(c)} The probabilities are close: $0.3968$ vs.\ $0.4096$ and
	$0.4290$ vs.\ $0.4096$. The reason: the sample (4 bonds) is small
	compared with the fund (50 bonds). Removing a few bonds hardly changes
	the share of junk bonds among the rest, so drawing without replacement
	is almost the same as drawing with replacement.

	The expected values are equal. For the hypergeometric distribution
	\[ \E[X]=\frac{nm}{N}=\frac{4\cdot10}{50}=0.8, \]
	and for the binomial distribution
	\[ \E[X]=np=4\cdot0.2=0.8. \]
	In both cases each picked bond is junk with probability $0.2$, so on
	average $4\cdot0.2=0.8$ of the 4 bonds are junk.

	\begin{answer}
		(a) hypergeometric with $N=50$, $m=10$, $n=4$;
		$P(X=0)=\frac{91390}{230300}\approx0.3968$,
		$P(X=1)=\frac{98800}{230300}\approx0.4290$ \quad
		(b) binomial with $(4,\,0.2)$; $P(X=0)=P(X=1)=0.4096$ \quad
		(c) close, since 4 is small compared with 50; $\E[X]=0.8$ in both cases
	\end{answer}
\end{solution}
